\documentclass[11pt]{article}
\usepackage[margin=1.12in]{geometry}
\usepackage{fontspec}
\setmainfont{Latin Modern Roman}
\usepackage[colorlinks=true,linkcolor=black,citecolor=black,urlcolor=blue]{hyperref}
\usepackage{microtype}
\setlength{\emergencystretch}{2em}
\usepackage{natbib}
\usepackage{fancyhdr}
\pagestyle{fancy}
\fancyhf{}
\renewcommand{\headrulewidth}{0pt}
\fancyfoot[L]{\footnotesize SLIPCASE · DENNING-2026-11-03 · 000\_\_RETURN\_PATH.txt}
\fancyfoot[R]{\thepage}
\title{Who Can Reopen the Decision?\\\large Customer Value and the Limits of Management Feedback}
\author{Research context: Watson Hartsoe\\Manuscript assembled by Codex}
\date{November 3, 2026 · Working Paper}
\begin{document}
\maketitle
\begin{abstract}
A company can learn a great deal from customers while making it nearly impossible for affected people to correct its decisions. This distinction narrows Steve Denning's argument that customer-centered firms have broken with the management model taught in business schools and sold by consultancies. Amazon's founding investment doctrine and McKinsey's historical work on customer value show that attention to customers fits comfortably inside a long-run financial strategy. Denning's earlier account of iterative delivery identifies a more consequential change: the movement of evidence into operating decisions. I distinguish that capacity from a right to reopen the decision itself. The distinction explains why faster iteration, better measurement, and outside verification may improve performance without repairing an accountability sink. It also suggests where to look for evidence that an institution has changed how it governs.
\end{abstract}
\noindent\textbf{Keywords} customer value, organizational learning, accountability, consulting, decision rights, management education

\section{What the strongest case would have to show}
Denning argues that business schools and strategy consultancies sustain one another through a shared language of hierarchy, financial optimization, and familiar analytical tools. Firms such as Amazon and Costco supposedly demonstrate the success of a different approach, centered on customers and sustained reinvestment \citep{denning2026businessschools}. The contrast carries an empirical burden. If the point is that a long investment horizon permits improvements that quarterly targets would kill, we need evidence about investment horizons. If the point is that customer benefit has independent weight, we need decisions in which the customer would be better off and investors were expected to be worse off. Evidence from cases where the two move together cannot distinguish the claims.

There is a further question about who can make the firm reconsider. An organization can adapt its offer repeatedly and still reserve to a small group the right to decide which consequences matter. This paper argues that responsiveness and the authority to demand review are distinct. The first can produce substantial benefits within an existing distribution of power. The second matters when that distribution filters out evidence against a decision that already has powerful sponsors. The argument is conceptual and comparative. The documents identify mechanisms and competing interpretations rather than their prevalence.

\section{The customer already had a place}
Amazon's 1997 shareholder letter joins customer focus, market leadership, and long-run shareholder value in a single chain. It treats near-term profitability as an unreliable guide to investments in leadership and treats future cash flows as decisive when they conflict with the appearance of accounting results \citep{amazon1997shareholderletter}. This is a clear example of a firm asking investors for patience. It is much weaker evidence that the customer displaced the investor as the ultimate beneficiary. That second claim requires evidence of a tradeoff where the investor case remained adverse even on management's own long horizon. A costly service launch is insufficient when management expected the expenditure to strengthen the franchise.

The claim of a historical break faces a different problem at McKinsey. Its 2000 account of economic value to the customer traces the approach to earlier staff work and describes comparing a proposed offer with what customers get from existing suppliers \citep{mckinsey2000delivering}. Such an estimate can guide an engineering investment. It can also inform a price increase that transfers part of the gain to the seller. A spreadsheet may contain a sophisticated representation of customer benefit while the decision it supports remains entirely financial. The presence of the customer in the model tells us little about which party keeps the improvement.

Porter's positive-sum account of differentiated rivalry supplies another check. Meeting the needs of different customer segments can expand an industry in his framework \citep{porter2008five}. Denning's more defensible criticism concerns the distance from that explanation to an operating routine. A theory of how differentiated offers improve industry outcomes leaves managers with the work of finding unmet needs and testing how to meet them. An argument about those routines has more purchase than a claim that earlier strategy had no conceptual room for better customer outcomes.

\section{The people a feedback loop can hear}
Denning's earlier description of Total Attorneys specifies an operating routine: cross-functional teams work in cycles, demonstrate the work to clients, and use what they hear to inform the next cycle \citep{denning2010radical}. This is a concrete mechanism for correction. It also gives us a way to see what correction can miss. In a legal-services workflow, the purchaser of a tool, the worker who uses it, and the person whose case goes through it can face different consequences. Consider a hypothetical feature that reduces the purchaser's handling time by making it more difficult for a claimant to correct a record. A team that learns primarily from the purchaser may correctly conclude that the feature is popular and become progressively worse at detecting the injury. Faster iteration compounds the omission when the omitted person has no way to introduce evidence.

Costco makes the distributional question visible without assuming a malfunction. The company describes how limited assortment, volume, and inventory turnover support a low-price operating model \citep{costco2025form10k}. Some customers can benefit from that choice while others incur a second trip to find an item the store declines to stock. Serving the latter group might compromise the low-price promise or might expose a needless constraint. A credible account of customer value needs a way to register the cost before adjudicating it. Completed purchases can inform the decision, but failed trips leave weaker records. The absence of those customers from the data may look like assent to the assortment.

\section{After the recommendation}
Davies identifies the accountability problem as a break between experienced consequences and the body responsible for the decision \citep{daviesTaming}. Graeber's account of the university employee who apologized when the carpenter failed to arrive gives the break a human form \citep{graeberRiseBullshitJobs}. The intermediary can provide real relief by acknowledging the problem. If the work ends there, the capacity decision survives unchanged. Whether the intermediary enables learning depends on where evidence of repeated failure goes and what authority waits at the other end.

This matters for consulting because outside involvement changes a decision before and after approval. Before approval, independent verification can make a disputed proposal more credible. Endorsement can also lend authority to a foregone conclusion. To distinguish them, look for the fate of an inconvenient finding. After approval, the relevant question is whether the client can revisit the assumptions, receive counterevidence, and revise the rule after the engagement ends. A recommendation might have been genuinely independent and still produce an accountability sink if the organization retains only a benchmark and no usable route for contesting it. Conversely, outside work can expand the organization's ability to revise its own decisions. Graeber's criterion of compelled pretense and Davies's feedback criterion thus diagnose different failures \citep{graeber2018bullshit,daviesTaming,denning2026businessschools}.

The same distinction changes what to ask of the business-school case method. Giving students the executive's chair can teach decision-making under uncertainty. If the exercise ends with an approved recommendation, students never confront the question of who can introduce consequential evidence after implementation. Including testimony at the outset leaves that authority question open. A second round could give someone subject to the decision a defined right to request review, compel a response, or seek a remedy. The instructional object would become the design of a decision process that can survive evidence against its author.

\section{An empirical test for the proposed divide}
A discriminating study would follow a single contested decision through four records: the original recommendation, the allocation or approval, the first material contrary evidence after implementation, and its disposition. At each stage the researcher would identify the actor entitled to introduce evidence and the actor authorized to change the outcome. A record of complaint closure would be coded separately from reconsideration. A successful exception for one complainant would be coded separately from a change to the decision rule. Independent review would count when the reviewer could reach an unwelcome conclusion and transmit it to an authorized decision-maker.

The same protocol could compare internal decisions with externally advised ones, and compare two decisions within the same firm. It would make a claim about management education testable only with additional evidence linking course design to how graduates actually structure review. The public documents examined here do not contain those records. They also cannot establish that selected public-company returns were caused by particular operating routines, or that those routines diffuse across firms. The argument survives those limits as a way to specify the evidence a stronger causal claim would require.

\section{Conclusion}
Denning's best insight is that management changes when organizations repeatedly learn from what their work does in the world. A strong account of that change needs to identify whose experience can interrupt the routine and who is authorized to respond. A purchaser can get better service, a consultant can deliver a sound analysis, and a manager can hit a customer outcome target while the people bearing a decision's costs remain unable to reopen it. The most telling evidence of changed governance would be a documented decision that an affected person could make the organization reconsider despite the approval, the metric, and the original sponsor.

\bibliographystyle{plainnat}
\bibliography{DENNING-2026-11-03__references}

\appendix
\section{Assembly instrument}
See the assembly receipt in \texttt{\_PROMPTS/} and the paper assembly appendix. The archive identifies the coverage and transcription limits of this record.
\section{Making history}
The researcher supplied the framing, source paste, and successive prompts. Codex selected the thesis, wrote the manuscript, and compiled the sources and lineage. The synthesis and proposed tests are interpretive. See the paper making history and source map for specific receipts.
\section{Replication path}
Open the standalone index, recover the embedded ZIP, and consult \texttt{000\_\_REBUILD.txt}. The manifest covers the archived inputs and derived payloads within the stated scope. The manuscript can be rebuilt with XeLaTeX and BibTeX. Research status is working paper. Human review is unverified.
\end{document}
